\section{Method}
\label{sec:method}

\begin{table}[t!]
\centering
\footnotesize
\caption{Seven distributional features. ``Entropy'' features are per-token NLL under Pythia-410M (\S\ref{sec:features}); surface features require no reference model. DDR analysis (\S\ref{sec:results-differential}) reveals entropy mean, TTR, and hapax ratio drift with depth while entropy max, min, and tail decay remain near-stationary; entropy std is intermediate.}
\label{tab:features}
\setlength{\tabcolsep}{4pt}
\begin{tabular}{@{}llp{3.0cm}@{}}
\toprule
\textbf{Feature} & \textbf{Source} & \textbf{Measures} \\
\midrule
\texttt{entropy\_mean} \textit{\scriptsize(NLL)} & Ref.\ NLL & Central tendency of token surprise \\
\texttt{entropy\_std} \textit{\scriptsize(NLL)}  & Ref.\ NLL & Spread of token surprise \\
\texttt{entropy\_max} \textit{\scriptsize(NLL)}  & Ref.\ NLL & Maximum token surprise \\
\texttt{entropy\_min} \textit{\scriptsize(NLL)}  & Ref.\ NLL & Minimum token surprise \\
\texttt{tail\_decay}   & Ref.\ NLL & Log-prob decay slope (ranks 100--1K) \\
\midrule
\texttt{ttr}           & Surface & Type-token ratio \\
\texttt{hapax\_ratio}  & Surface & Proportion of hapax legomena \\
\bottomrule
\end{tabular}
\end{table}

\subsection{Task Formulation}
\label{sec:task}

Given a text $x$, we predict its \emph{generation depth} $d \in \{0, 1, 2{+}\}$: human-written ($d{=}0$), first-generation LLM output ($d{=}1$), or text produced by two or more recursive LLM generations ($d{=}2{+}$).
A generation chain of depth $k$ is defined recursively: $x^{(0)} \sim \mathcal{C}_0$ (human corpus), $x^{(k)} = G(x^{(k-1)})$ for $k \geq 1$, where $G$ denotes stochastic autoregressive generation.
Depths beyond 2 are merged into $2{+}$ because feature convergence compresses inter-class margins at deeper levels (\S\ref{sec:results-differential}).

\subsection{Generation Chain Construction}
\label{sec:generation}

\looseness=-1 We construct generation chains using eight generators spanning five model families: \textbf{Pythia} (2.8B, 1B) \cite{pythia}, \textbf{OLMo} (OLMo-2-1B) \cite{olmo}, \textbf{Qwen} (Qwen3-1.7B, Qwen3-0.6B, Qwen3-8B-Base) \cite{qwen3}, \textbf{Llama} (Llama-3.2-3B) \cite{llama3}, and \textbf{Mistral} (Mistral-7B-v0.3) \cite{mistral7b}.
\looseness=-1 These span a 13$\times$ parameter range (0.6B--8B) across five architectures; Qwen's overrepresentation (3/8) is a limitation mitigated by leave-one-family-out (LOFO) evaluation (F1\,=\,0.686; \S\ref{sec:results-classification}).

\looseness=-1 For each generator, we sample approximately 1{,}000 seed passages from English Wikipedia as $\mathcal{C}_0$.
The same seed pool serves as gen-0 for all generators, controlling for input content; in LOGO evaluation this creates a potential train--test overlap for gen-0, addressed by our binary evaluation excluding gen-0 entirely (\S\ref{sec:results-baselines}).
Each seed is iteratively fed through the generator to produce a chain of depth 4 (i.e., gen-0 through gen-4), using the prompt ``\texttt{Continue the following text: \{text\}}''.
We generate up to depth~4 to empirically confirm convergence: 5-class LOGO yields macro F1\,$\approx$\,0.50, with depths 2--4 near-indistinguishable (gen-2/gen-3 per-class F1\,$<$\,0.27; Appendix~\ref{app:finer-depth}). This motivates the gen-2+ class collapse.
All generation uses temperature $\tau{=}0.7$, nucleus sampling with $p{=}0.9$, and maximum length of 256 new tokens.
After filtering chain breaks and passages too short for feature extraction, 37{,}058 texts remain across all depths (Appendix~\ref{app:dataset-stats}); depths 2--4 are collapsed into $2{+}$.

\subsection{Feature Extraction}
\label{sec:features}

\paragraph{Terminology.} Throughout this paper, we refer to per-token cross-entropy loss under a fixed reference model as ``entropy'' for conciseness; these are reference-model NLL values, not true information-theoretic entropy of the generating distribution.

We use interpretable distributional statistics so that the DDR partition remains transparent; end-to-end neural features improve classification ($+$10.0\,pp; Appendix~\ref{app:qlora}) but obscure the decomposition (LOFO confirms both feature groups contribute independently; Appendix~\ref{app:lofo}).

\looseness=-1 We extract seven distributional features per text (Table~\ref{tab:features}), covering the NLL distribution's location, spread, extremes, and tail structure, plus two reference-free surface diversity measures; six additional NLL statistics were rejected for redundancy ($|r| > 0.65$, $\Delta$F1\,$\leq$\,$-0.011$; Appendix~\ref{app:feature-candidate}).

\paragraph{Reference-model features.}
\looseness=-1 A held-out reference model (Pythia-410M, chosen to be smaller than all generators) computes per-token negative log-likelihood (NLL) for each text.
Cross-family reference validation (Appendix~\ref{app:cross-ref}) shows F1 spread of only 0.013 across Pythia, OLMo, and Qwen references, so DDR is reference-model agnostic; however, all tested references are causal LMs.
From the resulting NLL sequence, we extract five statistics: mean, standard deviation, maximum, and minimum of the token-level NLL, plus tail\_decay (the slope of the vocabulary-level averaged log-probability distribution in ranks 100--1000), which captures the power-law regime where depth-induced compression is most detectable.

\paragraph{Surface features.}
Two features capture lexical diversity without requiring a reference model: type-token ratio (TTR) and the proportion of words appearing exactly once (hapax\_ratio).

\paragraph{Feature complementarity.}
These features exhibit a \emph{differential depth response} (DDR), separating into drifting and non-drifting groups (\S\ref{sec:results-differential}; pairwise correlations in Appendix~\ref{sec:appendix-feature-correlation}).
TTR and hapax ratio are strongly correlated ($r = 0.97$); both are retained for marginal LOFO gains.

\subsection{Evaluation Protocol}
\label{sec:evaluation}

\paragraph{Cross-validation.}
\looseness=-1 We evaluate generalization via \textbf{LOGO} cross-validation: each fold holds out one generator entirely, so the classifier must generalize to a generator absent from training. This constitutes the strictest per-model protocol for assessing deployment against previously unseen generators.
LOGO does not isolate family-level effects; \textbf{LOFO} addresses this by holding out entire model families (binary F1\,=\,0.686; \S\ref{sec:results-classification}).
Generator subsets vary across experiments (Appendix~\ref{app:gen-subsets-section}).

\paragraph{Classifier.}
\looseness=-1 We use XGBoost \cite{xgboost} throughout ($n\_\text{estimators}{=}200$, $\max\_\text{depth}{=}6$, learning rate $0.1$; fixed without tuning), with identical hyperparameters across all methods in Table~\ref{tab:baselines}.
Alternative classifiers (CORAL \cite{coral_ordinal}, CORN \cite{corn_ordinal}, Frank-Hall \cite{frank_hall_ordinal}) yield comparable performance (Appendix~\ref{app:classifier}); hyperparameter sensitivity (max\_depth sweep) shows F1 variation $\leq$2.0\,pp, so discriminative power resides in the features, not the classifier.

\paragraph{Metrics and statistical testing.}
The primary metric is \textbf{macro-averaged F1} over the three depth classes, with 95\% bootstrap confidence intervals (CIs; $n{=}1{,}000$).
All $p$-values are two-sided with Bonferroni correction within each family of simultaneous tests.

\paragraph{Baselines.}
We compare the 7-feature ensemble against four single-score detection methods repurposed as depth predictors: Avg-LogProb, FastDetectGPT \cite{fast_detectgpt}, Binoculars \cite{binoculars}, and LD-Score \cite{ld_score}.
Each baseline score is used as a univariate feature in the same XGBoost framework under LOGO cross-validation to ensure a controlled comparison.

\paragraph{Length confound control.}
Since TTR and hapax\_ratio are mechanically length-dependent, we verify depth signal is not a length artifact via length-matched subsampling that equalizes length distributions across depth classes (\S\ref{sec:results-classification}; Appendix~\ref{app:length}).
